\subsection{Letter from Guenon to AKC 24VI1935}

\begin{quotex}
In 1935, \textbf{Rene Guenon} and \textbf{Ananda Coomaraswamy} began a correspondence that lasted 12 years. This letter shows their common mind, although their respective intellectual paths were quite different. An important point he makes regards the relationship between historical events and metaphysical principles; neither necessarily excludes the other. He also alludes what is called the ``Aryan Invasion Theory'' which brought the Primordial Tradition from the Arctic region to India in the form of Hinduism. To avoid confusion, keep in mind that Guenon emphasizes that the word ``Aryan'' does not designate a race.\footnote{The essay, ``The Darker Side of the Dawn'', can be found online.}
\end{quotex}


Cairo, 24 June 1935

Thank you for your letter and your shipment, which reached me at the same time.

There are some very interesting indications in ``The Darker Side of the Dawn'', which suggest many parallels with those found also in other traditions; it would be desirable if you get the opportunity to further develop these considerations later. As for your study of medieval aesthetics, the ideas you express fully coincide with what I think myself.

I am pleased to see that we are also in full agreement on the meaning of the words ``Aryan'' and ``non-Aryan'', which I could never consider as designating races.

As for the ``pre-Hindu'' traditions in India, without doubt, I have not sufficiently explained myself. It is well understood that all people are or have been in possession of traditions derived from a single source, but in a more or less distinct way. The Sumerian, Dravidian traditions, etc., seem to be based on forms most especially coming from certain secondary centers, whereas the ``Hindu'' tradition, which came from the north, is the one that comes most directly from the Primordial Tradition (for our Manvantara), shown everywhere as ``polar'' in origin. This of course has a direct link with the question of the ``Terrestrial Paradise'' which you alluded to, and which I mentioned in my book ``The King of the World''; this won't stop me from coming back to this topic again as you suggest. Concerning the analogy of historical events with the principles from which come their symbolic value (which does not exclude at all their factual reality), I often insisted on it; this is something that Westerners seem to have much difficulty in understanding in general.

I'm sorry to answer so briefly today but I hope now that our correspondence does not stay there.

Thank you again for what you tell me of my works and for making my works known in your circles.

Please accept, dear sir, my best wishes.


\flrightit{Posted on 2013-08-05 by Cologero}

\begin{center}* * *\end{center}

\begin{footnotesize}
\begin{sffamily}
\texttt{Avery Morrow on 2013-08-06 at 05:59 said:}

Here we are.

\url{http://www.scribd.com/doc/87506844/Coomaraswamy-The-Darker-Side-of-Dawn}


Whatever parallels Guenon was thinking of are beyond me on this early morning.

\hfill

\texttt{John Curley on 2013-08-06 at 06:23 said:}

Thank you for this link, Avery.


\hfill

\end{sffamily}
\end{footnotesize}

